\begin{topic}[id=ar_vr_xr,
             difficulty={Mittel},
             type={Überblick + Einordnung},
             prerequisites={Grundlagen Mensch-Maschine-Interaktion; Interesse an Sensorik und Visualisierung},
             practice={optional},
             owner={Dozententeam},
             updated={2026-04-09},
             tags={ AR, VR, XR, Haptik, Interaction, Tracking, Rendering, Ergonomie }]{AR/VR/XR \& haptische Displays}

\TopicDescription{XR-Systeme reichen von AR-Overlays bis zu immersiven VR-Umgebungen; haptische Displays erweitern den Feedbackkanal. Wir betrachten Interaktion, Tracking/SLAM, Latenzpfade, Rendering und Ergonomie. Ziel ist das Verständnis, wie Systemgrenzen (Motion-to-Photon, Field-of-View, Gewicht) Designentscheidungen prägen – und wo haptische Aktoren Mehrwert und Komplexität bringen.}

\TopicLeitfragen{\item Welche Latenz-/Tracking-Grenzen sind kritisch?
\item Wann lohnt haptisches Feedback?
\item Welche UX-Muster haben sich bewährt?
}
\TopicScope{\item Kein Game-Engine-Tutorial.
\item Keine Gerätevergleiche.
}
\TopicPractice{Optionale Praxisidee: UX-Skizze mit Latenzbudgets und haptischem Feedback für eine Aufgabe.}

\TopicKeywords{AR, VR, XR, Haptik, Interaction, Tracking, Rendering, Ergonomie}

\nocite{openxrSpec,xrHapticsSurvey}
\end{topic}
